\documentclass[11pt,a4paper]{article}

\usepackage[utf8]{inputenc}
\usepackage[T1]{fontenc}
\usepackage{amsmath,amsfonts,amssymb}
\usepackage{geometry}
\usepackage{enumitem}
\usepackage{titlesec}

\geometry{margin=2.5cm}

\titleformat{\section}{\large\bfseries}{\thesection.}{0.5em}{}
\titleformat{\subsection}{\normalsize\bfseries}{\thesubsection.}{0.5em}{}

\title{\textbf{Carnet de Suivi --- Projet de Market Making par Reinforcement Learning}}
\author{}
\date{}

\begin{document}

\maketitle

\section{Note méthodologique}

Le présent carnet de suivi a été initié tardivement dans le déroulement du projet.

Les premières phases de travail ont consisté en une exploration progressive des concepts (simulation de marché, définition d'états, premières implémentations), sans tenue systématique d'un journal formalisé.

Plutôt que de reconstruire artificiellement une chronologie précise, nous adoptons ici une démarche en deux temps :
\begin{itemize}
    \item une synthèse structurée des éléments déjà établis,
    \item puis un suivi rigoureux à partir de la phase actuelle du projet, centrée sur l'implémentation de l'algorithme PPO.
\end{itemize}

\textbf{Ce choix vise à privilégier la cohérence intellectuelle plutôt qu'une reconstruction chronologique artificielle.}

\section{Synthèse des travaux préliminaires}

\subsection{Objectif}

Le projet vise à développer une stratégie de market making à l'aide du reinforcement learning, dans un environnement simulé reproduisant certaines propriétés stylisées des marchés financiers.

\subsection{Modélisation}

Le problème est formulé comme un MDP :

\[
(S, A, P, R)
\]

avec :
\begin{itemize}
    \item un état basé sur des variables agrégées (micro-price, imbalance, RSI, inventaire),
    \item une action continue représentant les quotes bid/ask,
    \item une récompense fondée sur le PnL pénalisé par le risque d'inventaire.
\end{itemize}

\subsection{Simulateur}

Un environnement simplifié a été développé :
\begin{itemize}
    \item dynamique du mid-price et du spread,
    \item probabilités de fill dépendant des quotes,
    \item gestion de l'inventaire.
\end{itemize}

\textbf{Remarque critique :}  
Cette phase a principalement consisté en une construction expérimentale visant à rendre le problème manipulable, au prix d'hypothèses simplificatrices fortes.

\section{Début du suivi formalisé}

\subsection*{Entrée du 13 avril 2026 --- Définition de l'architecture PPO}

\textbf{Objectif :}  
Initier l'implémentation d'une méthode Actor-Critic basée sur PPO.

\textbf{Décisions prises :}
\begin{itemize}
    \item Adoption d'une politique continue paramétrée par un réseau de neurones
    \item Utilisation d'une distribution gaussienne suivie d'une transformation $\tanh$
    \item Séparation Actor / Critic
\end{itemize}

\textbf{Analyse :}  
Lors du passage a une politique d'actions continue, le papier n'indiquait pas de quelle manière se définissait cette loi. 
Nous avons alors choisi une politique suivant une loi log-normale pour conserver la positivité des actions et une certaine simplicité. 

\subsection*{Entrée du 14 avril 2026 --- Implémentation de l'Actor}

\textbf{Objectif :}  
Structurer le réseau Actor (\texttt{ActorNet}) et définir la fonction de sampling.

\textbf{Travail réalisé :}
\begin{itemize}
    \item Définition d'un réseau fully-connected avec activations \texttt{tanh}
    \item Paramétrisation d'une loi gaussienne sur un espace latent
    \item Transformation via $\tanh$ pour contraindre l'action
\end{itemize}

\textbf{Question centrale :}  
Quel est le rôle exact de la fonction \texttt{sample\_action} ?

\textbf{Clarification :}
Cette fonction doit :
\begin{itemize}
    \item échantillonner une action selon la politique $\pi_\theta$
    \item retourner la log-probabilité associée
    \item éventuellement retourner l'action non transformée (pré-$\tanh$)
\end{itemize}

\textbf{Importance :}  
Ces éléments sont nécessaires pour calculer le ratio PPO :

\[
r_t(\theta) = \frac{\pi_\theta(a_t|s_t)}{\pi_{\theta_{old}}(a_t|s_t)}
\]

\subsection*{Entrée du 14 Avril 2026 --- Définition de la perte de l'acteur}
On aborde la définition de la fonction de perte associée a la politique. Contrairement au critic, la cible n'est pas clairement identifiée, car il n'existe pas de "bonne action" connue a priori. 

On introduit alors un signal d'apprentissage, 
\textbf{L'Avantage} calculé comme suit
\[
A_t = G_t - V(s_t)
\]
où 
\begin{itemize}
    \item $G_t$ représente le gain observé ici
    \item $V(s_t)$ correspond au gain attendu sous la politique actuelle
\end{itemize}
L'avantage mesure donc a posteriori la qualité d'une action, et nous nous servons de ce signal pour ajuster la politique
\[
\log\pi_\theta (a_t|s_t)\cdot A_t
\]
Ainsi, les actions ayant un avantage positif sont renforcées, tandis que celles ayant un avantage négatif sont pénalisées.

On définit alors la perte de l'acteur comme 
\[
L_{actor} = -\mathbb{E}[\log\pi_\theta (a_t|s_t)\cdot A_t]
\]
Cette perte est implémentée dans le module Loss. 

\section*{Entrée du 15 avril 2026 : Premiere implémentation de la loss d'acteur, implémentation de la boucle d'apprentissage}
La perte est calculée comme une simple moyenne pondérée avec les avantages, sans bonus d'entropie. 

On met en place une première boucle d'apprentissage locale pour l'acteur, réévalue la politique courante sur les actions et entraine sur la perte précédemment définie. 
\end{document}